\appendix

\section{Dominant morphisms}
\label{sec:dom}

Dominant morphisms can be defined in a variety of different but equivalent ways.

\begin{theorem}[{\cite[Section AG, Theorem 17.3]{borel91}}]\label{borel91}
	Let $X \subset \mathbb{C}^m$ be an algebraic set and $Y \subset \mathbb{C}^n$ be a variety.
	Then the following are equivalent for any morphism $f :X \rightarrow Y$:
	\begin{enumerate}
		\item $f$ is dominant.
		\item For some irreducible component $X'$ of $X$, there exists a point $x \in X'$ such that $x$ is a non-singular point of $X'$ and $\rank \diff f(x) = \dim Y$.
		\item There exists a Zariski open subset $U \subset X$ where for each $x \in U$, $x$ is a non-singular point of $X$ and $\rank \diff f(x) = \dim Y$.
	\end{enumerate}
\end{theorem}

It follows immediately from \Cref{borel91} that for any varieties $X,Y$,
the existence of a dominant morphism from $X$ to $Y$ implies $\dim X \geq \dim Y$.
As can be seen by \Cref{t:deg},
more powerful statements can be attained relating to $f$ if $\dim X=\dim Y$.
To prove \Cref{t:deg}, we require the following two results.

\begin{corollary}[{\cite[Section 8, Corollary 1]{mumford}}]\label{mumford}
	Let $X \subset \mathbb{C}^m$ and $Y \subset \mathbb{C}^n$ be varieties and $f:X \rightarrow Y$ be a dominant morphism.
	Then there exists a non-empty Zariski open subset $U \subset Y$ such that $U \subset f(X)$, and for every $y \in U$, every irreducible component of the algebraic set $f^{-1}(y)$ has dimension $\dim X - \dim Y$.
\end{corollary}

\begin{corollary}[{\cite[Proposition 7.16]{Harris}}]\label{c:harris}
	Let $X \subset \mathbb{C}^m$ and $Y \subset \mathbb{C}^n$ be varieties and $f:X \rightarrow Y$ be a dominant morphism.
	Suppose that there exists a non-empty Zariski open subset $U \subset Y$ where $|f^{-1}(y)|<\infty$ for every $y \in Y$.
	Then there exists a $k \in \mathbb{N}$ and a non-empty Zariski open subset $U' \subset U$ where $|f^{-1}(y)| = k$ for every $y \in U'$.
\end{corollary}

\begin{proof}[Proof of \Cref{t:deg}]
	It is immediate that \ref{t:deg3} implies \ref{t:deg2}.
	Fix $U \subset Y$ to be the Zariski open set from \Cref{mumford}.
	An algebraic set is a finite set if and only if it is zero dimensional.
	Hence, \ref{t:deg1} and \ref{t:deg2} are equivalent.
	Suppose that $\dim X=\dim Y$.
	Fix $U' \subset U$ to be the non-empty Zariski open subset from \Cref{c:harris}.
	Define the set
	\begin{align*}
		C := \{x \in X : x \text{ is singular or } \rank \diff f(x) <n\}.
	\end{align*}
	By \Cref{borel91},
	$C$ is a proper Zariski closed subset of $X$.
	As $\dim C < \dim X = \dim Y$,
	the set $f(C)$ is not dense in $Y$.
	It now follows that the complement of the Zariski closure of $f(C)$ in $U$ is a non-empty Zariski open subset of $Y$.
	Hence, \ref{t:deg1} implies \ref{t:deg3},
	concluding the proof.
\end{proof}







\section{Proof of \texorpdfstring{\Cref{l:xgd}}{Lemma}}
\label{sec:xgd}

We begin with the following technical lemma describing when we can ``rotate'' a framework into our set $X_{G,d}$.

\begin{lemma}\label{l:rotate}
	Let $G=(V,E)$ be a graph with $|V| \geq d+1$,
	and fix a sequence of $d$ vertices $v_1,\ldots, v_d$.
	For each realisation $p \in (\mathbb{C}^d)^V$,
	define the $(d-1) \times (d-1)$ symmetric matrix
	\begin{align*}
		\mathbb{G}(p) := 		
		\begin{bmatrix}
			(p_{v_2} - p_{v_1})^T \\
			\vdots \\
			(p_{v_d} - p_{v_1})^T
		\end{bmatrix}
		\begin{bmatrix}
			p_{v_2} - p_{v_1} & \cdots & p_{v_d} - p_{v_1}
		\end{bmatrix}.
	\end{align*}
	Then there exists a realisation $q \in X_{G,d}$ congruent to $p$ if $\mathbb{G}(p)$ only has non-zero leading principal minors.\footnote{A \emph{leading principal minor (of order $n$)} of a square matrix is the determinant of the matrix formed by taking the first $n$ rows and columns.}
\end{lemma}

The conditions stated in \Cref{l:rotate} seem rather bizarre,
especially since no such conditions are required if we restrict ourselves to real realisations.
It is also rather easy to see they are not necessary either:
simply choose any realisation $p \in X_{G,d}$ where $p_{v_i} = \mathbf{0}$ for each $i \in \{1,\ldots,d\}$.
To see why we need to be so cautious,
take $G=(V,E)$ to be the complete graph with 4 vertices $v_1,v_2,v_3,v_4$,
and let $p$ be the 3-dimensional realisation of $G$ where
\begin{align*}
	p_{v_1} = (0,0,0), \qquad p_{v_2} = (1,0,0), \qquad p_{v_3} = (2,1,i), \qquad p_{v_4} = (0,1,0).
\end{align*}
In this specific case, there are no realisations in $X_{G,d}$ that are congruent to $p$.
The framework $(G,p)$ also has the interesting properties that its vertices affinely span $\mathbb{C}^3$ and $\|p_v-p_w\|^2 >0$ for every edge $vw \in E$.
Since the matrix
\begin{align*}
	\mathbb{G}(p) =
	\begin{bmatrix}
		1 & 0 & 0 \\
		2 & 1 & i
	\end{bmatrix}
	\begin{bmatrix}
		1 & 2 \\
		0 & 1 \\
		0 & i
	\end{bmatrix}
	=
	\begin{bmatrix}
		1 & 2 \\
		2 & 4
	\end{bmatrix}
\end{align*}
has rank 1,
our constructed realisation is not a counter-example to \Cref{l:rotate}.

\begin{proof}
	Fix $p \in (\mathbb{C}^d)^V$ with the stipulated properties.
	With this,
	define $p^1 \in (\mathbb{C}^d)^V$ to be the realisation where $p^1_v = p_v - p_{v_1}$ for each $v \in V$.
	If $d=1$ then $p^1 \in X_{G,1}$,
	so suppose that $d \geq 2$.
	We now observe that the matrix $\mathbb{G}(p^1) = \mathbb{G}(p)$.
	In fact, a much stronger statement is true:
	if $q$ is a congruent realisation then $\mathbb{G}(q) = \mathbb{G}(p)$.
	
	We now form the sequence of congruent realisations $p^1,\ldots,p^d$ in the following inductive way.
	Fix $n \in \{2,\ldots,d\}$,
	and suppose that $p^1,\ldots,p^{n-1}$ have already been constructed such that $[p^{n-1}_{v_k}]_j =0$ if $k \leq \min \{j, n-1\}$.
	For each $k \in \{1,\ldots,d\}$,
	fix $e_k$ to be the vector with $[e_k]_j = 1$ if $j=k$ and $[e_k]_j = 0$ otherwise.
	We observe here that the linear space spanned by the vectors $e_1,\ldots, e_{n-2}$ contains the points $p^{n-1}_{v_2},\ldots, p^{n-1}_{v_{n-1}}$.
	Fix $z$ to be the vector formed from $p^{n-1}_{v_{n}}$ by replacing its first $n-2$ coordinates with zeroes.
	It is immediate that $z \cdot e_j = 0$ for all $j \in \{1 , \ldots, n-2\}$.
	
	Suppose for contradiction that $\|z\|^2 = 0$,
	and fix $y = p^{n-1}_{v_{n}} - z$.
	Set $A = [p^{n-1}_{v_2} ~ \cdots ~ p^{n-1}_{v_{n-1}}]$.
	The determinant of $A^T A$ is the leading principal minor of order $n-2$ of $\mathbb{G}(p) = \mathbb{G}(p^{n-1})$,
	and so is non-zero.
	Hence, the column span of $A$ is exactly the linear space spanned by $e_1,\ldots, e_{n-2}$;
	importantly,
	this implies that $y$ is contained in the column span of $A$.
	The leading principal minor of order $n-1$ of $\mathbb{G}(p) = \mathbb{G}(p^{n-1})$ is the determinant of the $(n-1) \times (n-1)$ matrix
	\begin{align*}
		B &=		
		\begin{bmatrix}
			(p^{n-1}_{v_2})^T \\
			\vdots \\
			(p^{n-1}_{v_{n}})^T
		\end{bmatrix}
		\begin{bmatrix}
			p^{n-1}_{v_2} & \cdots & p^{n-1}_{v_{n}}
		\end{bmatrix}\\
		&=
		\begin{bmatrix}
			A^T A & A^T (y+z) \\
			(y+z)^T A & (y+z)^T (y+z)
		\end{bmatrix}\\
		&=
		\begin{bmatrix}
			A^T A & A^T y \\
			y^T A & y^T y
		\end{bmatrix}\\
		&= 
		\begin{bmatrix}
			A^T \\
			y^T
		\end{bmatrix}
		\begin{bmatrix}
			A & y
		\end{bmatrix}.
	\end{align*}
	As $y$ is contained in the column span of $A$,
	the matrix $[A ~ y]$ has rank at most $n-2$.
	However, this implies $\det B = 0$,
	contradicting our leading principal minor assumption for $\mathbb{G}(p)$.
	Hence, $\|z\|^2 \neq 0$.
	
	Fix $x := z/\|z\|^2$.
	Define the $d \times d$ matrix $\tilde{M}$ where $\tilde{M}^T = [ e_1 ~ \cdots ~ e_{n-2} ~ x ~ \mathbf{0} ~ \cdots ~ \mathbf{0}]$.
	We note that the matrix $\tilde{M}$ is an isometry (in the quadratic form sense) from the linear subspace of $\mathbb{C}^d$ spanned by $p^{n-1}_{v_2},\ldots, p^{n-1}_{v_{n}}$ to the linear subspace $\mathbb{C}^{n-1} \times \{0\}^{d-n+1}$.
	By Witt's theorem (see, for example,~\cite[Theorem 11.15]{roman}),
	there exists a matrix $M \in O(d, \mathbb{C})$ where $Mp^{n-1}_{v_j} = \tilde{M}p^{n-1}_{v_j}$ for each $j \in \{1,\ldots,n\}$.
	With this we now fix $p^n_v = Mp^{n-1}_v$ for each $v \in V$.
	The result now follows from completing the inductive argument and fixing $q = p^d$.
\end{proof}


With this, we can prove \Cref{l:xgd}.


\begin{proof}[Proof of \Cref{l:xgd}]
	For every $p \in (\mathbb{C}^d)^V$,
	let $\mathbb{G}(p)$ be the matrix defined in the statement of \Cref{l:rotate}.
	With this, we fix the proper algebraic set
	\begin{align*}
		Z := \left\{ p \in (\mathbb{C}^d)^V :
			\begin{array}{l}
				\text{$\mathbb{G}(p)$ has a zero leading principal minor,}\\
				\text{or $p$ does not affinely span $\mathbb{C}^d$}
			\end{array}
		 \right\}.
	\end{align*}
	Since $Z$ contains hypersurfaces in $(\mathbb{C}^d)^V$ and $Z \neq (\mathbb{C}^d)^V$
	it has dimension $d|V|-1$.
	By \Cref{l:rotate},
	the set $f_{G,d}((\mathbb{C}^d)^V \setminus Z)$ is a subset of the image of $\tilde{f}_{G,d}$.	
	Hence, the image of $\tilde{f}_{G,d}$ is Zariski dense in the image of $f_{G,d}$.
	
	Since $\dim \ell_d(G) = \rank \diff f_{G,d}(p)$ for any general realisation $p\in (\mathbb{C}^d)^V$,
	it follows from \Cref{t:asimowroth} that $G$ is $d$-rigid if and only if $\dim \ell_d(G) = d|V|- \binom{d+1}{2}$.
	Suppose that $G$ is not $d$-rigid.
	The extension of \Cref{t:asimowroth} to complex realisations gives that the set $C_d(G,p)$ contains infinitely many points for almost all realisations $p$;
	in particular,
	it can be used to prove that for almost all $p \in (\mathbb{C}^d)^V \setminus Z$,
	there exists a sequence $(p^n)_{n \in \mathbb{N}}$ of realisations where $p^n$ is not congruent to either $p^m$ or $p$ for each $m \neq n$, and $p^n \rightarrow p$ as $n \rightarrow \infty$ in the standard complex norm for $(\mathbb{C}^d)^V$ (note that this is an actual metric and not the quadratic form $\| \cdot\|^2$).
	An application of \Cref{l:rotate} to both $p$ and each $p^n$ gives that $\tilde{f}_{G,d}^{-1}\left(f_{G,d}(p)\right)$ also contains infinitely many points as is required.
	
	We now suppose that $G$ is $d$-rigid,
	i.\,e., $\dim \ell_d(G) = d|V|- \binom{d+1}{2}$.
	We split the proof into three cases.
	
	\textbf{(Case 1: $f_{G,d}(Z)$ is Zariski dense in $\ell_d(G)$ and $|V| \geq d+2$.)}
	In this case,
	the restriction of the map $f_{G,d}$ to $Z$ and $\ell_d(G)$, which we now denote by $g : Z \rightarrow \ell_d(G)$, is dominant.
	Hence, by \Cref{borel91},
	there exists a non-empty Zariski open subset $U \subset Z$ where 
	\begin{align*}
		\rank \diff g(p) = \dim \ell_d(G) = d|V|- \binom{d+1}{2}
	\end{align*}
	for each $p \in U$.
	Since $\rank \diff g(p) \leq \rank \diff f_{G,d}(p)$,
	we have that $\rank \diff f_{G,d}(p) = d|V| - \binom{d+1}{2}$ for each $p \in U$.
	As $Z$ has dimension $d|V|-1$,
	it contains an irreducible component of dimension $d|V|-1$.
	It follows that the nullity of $\diff g(p)$ is $\binom{d+1}{2} - 1$ for each $p \in U$.	
	If $q$ is congruent to $p \in Z$ then $\mathbb{G}(q) = \mathbb{G}(p)$ and $q$ is affinely spanning if and only if $p$ is,
	hence $Z$ is closed under congruences.
	This implies that for each $p \in U$,
	the map from the affine congruences of $\mathbb{C}^d$ to the corresponding congruent realisation of $p$ in $Z$ is not injective,
	as otherwise the nullity of $\diff g(p)$ would be at least $\binom{d+1}{2}$ for each $p \in U$.
	Hence, any realisation in $U$ cannot affinely span $\mathbb{C}^d$.
	The set of realisations of $G$ that do not affinely span $\mathbb{C}^d$ is the intersection of $|V| - d \geq 2$ pairwise-distinct hypersurfaces,
	and so $U$ is contained in a $(d|V|-2)$-dimensional subvariety in $Z$.
	However, this now contradicts that $U$ is a non-empty Zariski open subset of the $(d|V|-1)$-dimensional algebraic set $Z$.
	
	
	\textbf{(Case 2: $f_{G,d}(Z)$ is not Zariski dense in $\ell_d(G)$ and $|V| \geq d+2$.)}
	Fix the non-empty Zariski open set
	\begin{align*}
		X := \left\{ p \in X_{G,d} : \tilde{f}_{G,d}(p) \text{ is not contained in the closure of $f_{G,d}(Z)$ }   \right\}.
	\end{align*}
	Choose any $p \in X$.
	By \Cref{l:rotate}, for each realisation $p' \in (\mathbb{C}^d)^V$ that is equivalent to $p$,
	there exists some other realisation $p'' \in X_{G,d}$ that is congruent to $p'$.
	Hence, it now suffices to prove that $p$ is congruent to exactly $2^d$ realisations (including itself) in $X_{G,d}$.
	
	Choose any $q \in X_{G,d}$ congruent to $p$,
	and let $A \in O(d,\mathbb{C})$ be the matrix that maps the vertices of $p$ to the vertices of $q$.
	As $p,q \in X_{G,d}$, we have 
	\begin{align*}
		\sum_{n=1}^k A_{j,n} [p_{v_{k+1}}]_n = [q_{v_{k+1}}]_j = 0
	\end{align*}
	for each $1 \leq k < j \leq d$,
	hence $A_{j,k}=0$ for each $1 \leq k < j \leq d$.
	Since $A^T$ maps the vertices of $q$ to the vertices of $p$,
	we similarly have $A_{j,k}=0$ for each $1 \leq j < k \leq d$.
	As the vertices of $p$ linearly span $\mathbb{C}^d$,
	the set $\tilde{f}_{G,d}^{-1}\left(f_{G,d}(p) \right)$ is in one-to-one correspondence with the diagonal orthogonal matrices.
	The set of diagonal orthogonal matrices is exactly the set of diagonal matrices $A$ where $A_{j,j} = \pm 1$ for each $j \in \{1,\ldots, d\}$,
	hence there are $2^d$ of them.
	This now concludes the proof.
	
	\textbf{(Case 3: $|V| =d+1$.)}	
	Since $G$ is $d$-rigid, it can be easily verified that $G$ must be the complete graph on $d+1$ vertices.
	Furthermore,
	every element $p \in (\mathbb{C}^d)^V \setminus Z$ is only equivalent to congruent realisations (\cite[Corollary 8]{Gortler2014}),
	and each equivalent realisation is also contained in $(\mathbb{C}^d)^V \setminus Z$.
	We now define $X = X_{G,d} \setminus Z$ and repeat the same technique as was utilised in Case~2.
	This now concludes the proof.
\end{proof}





